% ID: SEL-D0042-OB-E03
% ESTADO: cotejada. Segunda pasada visual y borrador compilado revisados.
% Fuente: Selectividad/sel_2014_matematicas_aplicadas/EXAMEN -ANDAL-1.pdf
% SHA-256: 375424f4e0d3c0409619ebaae4ae4d20c441a2585efe3dc74eb98d765114311e
% Página: 2; opción B; ejercicio original 3.

Un estudio estadístico de la producción de una fábrica de batidoras determina que el 4.5\% de las batidoras presenta defectos eléctricos, el 3.5\% presenta defectos mecánicos y el 1\% presenta ambos defectos. Se escoge al azar una batidora.
\begin{enumerate}[label=\alph*)]
\item \textbf{(1 punto)} Calcule la probabilidad de que no tenga ninguno de los dos defectos.
\item \textbf{(1 punto)} Calcule la probabilidad de que tenga un defecto mecánico sabiendo que tiene un defecto eléctrico.
\item \textbf{(0.5 puntos)} Justifique si los sucesos ``tener un defecto eléctrico'' y ``tener un defecto mecánico'' son independientes. \textquestiondown{}Son incompatibles?
\end{enumerate}
